% s-055-p055-m002; original PDF 55, 56, 57, 58, 59
\Aphorism{s-055-p055-m002}{55}
\Sperr{Etwas} \Sperr{vom} \Sperr{Altruismus.} Wann empfinden wir die „Gesellschaft“ als liebenswürdig, als „Object möglicher Liebe“? Wann und wie lieben wir überhaupt, wir Egoisten? Nun, im Allgemeinen sicherlich aus der Ferne, mit hinlänglichem Spielraum für unser Ego, das „idealisiren“, umformen, zurechtmachen, fälschen, aneignen, assimiliren will; wir lieben nicht die Dinge, wie sie sind, sondern unsere Vorstellung von ihnen, irgend eine uns angepasste Bezeichnung und Symbolik der Dinge, an der sich unsere Willkür, unser Eigenthumstrieb, unsere „Liebe“ ungehindert entladen kann. Sie selbst, die Dinge an sich, sind viel zu hart, fremd, objectiv, um in dieser animalischen Weise geliebt zu werden. Und wie lieben wir die Gesellschaft? Nach meinen Erfahrungen ebenfalls aus der Ferne, sinnlich und geistig genommen, aus jener „blauen“ Ferne, in der alle hässlichen, schrillen, gezackten Einzelheiten gedämpft und überhaucht erscheinen und auch das wimmelnde Collectivum zu etwas Rundem, Einheitlichem, physiologisch Ganzem zusammenrinnt. Eine Möglichkeit, die Vorstellung „Gesellschaft“ mit Lust- oder Unlustgefühlen zu betonen, tritt für mich genau auf dem Punkte ein, wo diese Gesellschaft als leibliches und abgeschlossenes Object, gewissermassen als colossales Individuum, zu meinen Sinnen redet, wo die rein begriffliche Addition, durch die jedes Collectivum als Summe seiner Einzelnen, und nicht mehr, definirt ist, einmal in den Hintergrund tritt und sich dafür der Begriff eines Organismus einschiebt. Ich bin ausser Stande, mit der Bereitwilligkeit, die ein heutiger Socialethiker fordert, die Gesellschaft \Sperr{durchgängig} als organisirtes Wesen zu empfinden; aber ich bin \Sperr{zeitweilig} dazu im Stande. Und immer, wenn irgend eine verliebte Zärtlichkeit für „Andere“, eine weichherzige altruistische Anwandlung mit kribbeligen Schmetterlingsfüssen über meine Seele hinläuft, ist es sicherlich die physiologische Realität, nicht das „moralische Ideal“, in die ich mich einen Augenblick verliebe. Irgend ein Sinn bewirkt eine confocale Vereinigung der Reizstrahlen und dadurch eine regelrechte Gefühlserwärmung: je besser die synthetische Zusammenschmelzung der gesellschaftlichen Vielheitskategorie zur animalen Einheitskategorie gelingt, desto leichter fängt ein wenig Altruismus Feuer. Wenn ich inmitten eines lärmenden Häufleins Menschlichkeit die einzelnen Stimmen höre und die einzelnen Gebärden sehe, so fehlt mir jede Sympathie, jedes Gefühlscorrelat für eine angebliche Resultante aller dieser zersplitterten Componenten; ganz anders, wenn dieses selbe Stimmengewirr, weit fort von mir, als unlösliche akustische Verbindung ins Schweigen der Ferne versinkt, immer leiser, voller, ausgerundeter, liebenswürdiger \Dots{4} Ebenso kann mich der Anblick zahlloser Fussspuren auf weichem Boden oder im Schnee ganz zart altruistisch stimmen, gleichermassen der aus Ballsälen und Theatern bekannte Duft von Parfums, Glacéhandschuhen und menschlicher Ausdünstung — jener Duft, der eine Gesellschaft aus den „höheren Ständen“ charakterisirt und gewissermassen zum ausdünstenden \Foreign{animal} vereinheitlicht. Auch die dem Religiösen angenäherte Empfindung, die Einen inmitten einer jubelnden oder singenden Volksmenge überkommt, gehört hierher. Ja, auf die Gefahr hin, in eine haarscharfe Spitze des Raffinements auszulaufen, möchte ich behaupten, dass selbst im \Sperr{Lautwerth} eine Erleichterung oder Erschwerung zur sympathischen Betonung der Begriffe und Vorstellungen liegen kann; ich würde mich z. B. eher für eine „Gesellschaft“ als für einen „Staatsverband“ begeistern: die Gesellschaft, mit der doppelten Liquida und den beiden Sibilantes, ist — ja, wie nenne ich es doch? feuchter, sinnlicher, lebendiger, physiologischer, mehr zum Liebkosen \Dots{4} Die „Gemeinde“ oder der „Staat“, sachlich ebenso abstract und zusammenaddirt wie jene, hat nicht die Geschicklichkeit, phonetisch das Gegentheil zu scheinen!
\par

Diese altruistischen Gefühlsverzweigungen, die ein Anderer durch entsprechende Singularitäten seiner eigenen Erfahrung ersetzen wird, mögen spitz, seltsam, subjectiv und Schlimmeres sein: jedenfalls geben sie uns, den individualistischen Menschen, einen Genuss am Gegentheil unserer Moral! Ich vermuthe, dass dieser Genuss den gewohnheitsmässigen altruistischen Dickhäutern, die ihre Menschenliebe in grossen Alltagsportionen zu sich nehmen, entgehen wird. Die Gewohnheit stumpft gegen Nüancen ab, nur die Ausnahme schmeckt und wird geschmeckt; einer Sache fernbleiben ist die einzige Möglichkeit, sich eine feine Zunge dafür zu bewahren, heisse diese Sache nun Kunst, Religion, Moral, Geliebte, Landschaft oder sonstwie. Sicherlich werden die zartesten Entzückungen und Wohlgerüche des moralischen Lebens nicht von den professionellen Tugendbolden, sondern von den Bösen und Halbmoralischen genossen. Wir als Individualisten haben längst auf die wundervollen Reize der Selbstliebe verzichtet; dafür schenkt uns unsere Feindin, die Socialethik, gelegentlich einen Augenblick im Paradiese. Diesen berauben die durchgängig altruistischen Menschen seiner Kraft und Süssigkeit, indem sie ihn zur ganzen Lebenslinie in die Länge ziehen und die Ausnahme zur Regel (zur „Pflicht“!) verdünnen; vielleicht hält sie der Egoismus, der ihnen eine gefährliche und ketzerische Sache ist, durch \Sperr{seine} paradiesischen Augenblicke dafür schadlos. Aber trotz dieser gegenseitigen Entfremdung im Verstehen und Geniessen fehlt es nicht an Gründen, das was wir in jenen altruistischen Moment-Anwandlungen gelernt haben, mit zur Psychologie des altruistischen Normalzustandes heranzuziehen. Ohne dies des Näheren auszuführen, können wir uns mit der These begnügen, dass die Hingebung an ideale Collectiva, so abstract und moralisch-vernünftig sie sich auch verkleiden mag, doch zuletzt zu den \Sperr{sinnlichen} Emotionen zu rechnen ist, insofern die begriffliche Construction „Gesellschaft“ sich zu einer Art physiologischer Leiblichkeit erwärmt und belebt haben muss, um sich des Weiteren gefühlsmässig färben und betonen zu lassen. Erst wenn man die Vielheit nicht nur als Einheit denken, sondern (in irgend einer directen oder reflectirten Form) als Einheit sehen, hören, riechen, schmecken, tasten kann, erst dann bemächtigt sich ihrer die moralische Theilnahme, als geistige Ausstrahlung einer im tiefsten Grunde sinnlichen, im allertiefsten geschlechtlichen Energie.
\par

